% Part: turing-machines
% Chapter: undecidability
% Section: enumerating-tms

\documentclass[../../../include/open-logic-section]{subfiles}

\begin{document}

\olfileid[tr]{tur}{und}{enu}
\olsection{Turing Makinelerinin Numaralandırılması}

\begin{explain}
Bütün Turing makineleri kümesinin !!{enumerable} olduğunu gösterebiliriz.
Bu, her Turing makinesinin sonlu biçimde betimlenebilmesinden çıkar. Durumlar
kümesiyle şerit söz dağarcığı sonlu kümelerdir. Geçiş fonksiyonu
$Q\times\Sigma$'dan
$Q\times\Sigma\times\{\TMleft,\TMright,\TMstay\}$'ye bir kısmi
fonksiyondur; dolayısıyla o da tanımlı olduğu sonlu sayıdaki argüman çifti
için değerleri listelenerek belirtilebilir.

Buraya kadarı doğrudur, fakat ince bir ayrım vardır. Turing makinelerinin
tanımı, durumlar kümesiyle şerit alfabesinin hangi nesneleri !!{element}
olarak içerebileceğine hiçbir kısıtlama getirmedi. Dolayısıyla örneğin her
gerçel sayı için, teknik olarak, o sayıyı durum olarak kullanan bir Turing
makinesi vardır. Ancak Turing makinesinin \emph{davranışı}, hangi nesnelerin
durum ve söz dağarcığı olarak kullanıldığından bağımsızdır.
\olref{fig:variants} içindeki iki Turing makinesini ele alın.
\begin{figure}
\begin{center}
\begin{tikzpicture}[->,>=stealth',shorten >=1pt,auto,node distance=2.8cm,
                    semithick]
  \tikzstyle{every state}=[fill=none,draw=black,text=black]

  \node[initial,state]         (A)              {$q_0$};
  \node[state]         (B) [right of=A] {$q_1$};

  \path (A) edge [bend left] node {\TMtrans{\TMstroke}{\TMstroke}{\TMright}} (B)
        (B) edge [loop above] node {\TMtrans{\TMblank}{\TMblank}{\TMright}} (B)
            edge [bend left] node {\TMtrans{\TMstroke}{\TMstroke}{\TMright}} (A);
\end{tikzpicture}\\
\begin{tikzpicture}[->,>=stealth',shorten >=1pt,auto,node distance=2.8cm,
  semithick]
\tikzstyle{every state}=[fill=none,draw=black,text=black]

\node[initial,state]         (A)              {$s$};
\node[state]         (B) [right of=A] {$h$};

\path (A) edge [bend left] node {\TMtrans{A}{A}{\TMright}} (B)
(B) edge [loop above] node {\TMtrans{\TMblank}{\TMblank}{\TMright}} (B)
edge [bend left] node {\TMtrans{A}{A}{\TMright}} (A);
\end{tikzpicture}
\end{center}
\caption{\emph{Çift Sayı} makinesinin çeşitlemeleri}
\ollabel{fig:variants}
\end{figure}
Bu iki diyagram iki makineye karşılık gelir: Şerit alfabesi
$\Sigma=\{\TMendtape,\TMblank,\TMstroke\}$ ve durumlar kümesi
$\{q_0,q_1\}$ olan $M$ ile alfabesi
$\Sigma'=\{\TMendtape,\TMblank,A\}$ ve durumları $\{s,h\}$ olan $M'$.
Ancak bunun dışında yönergeleri aynıdır: $M$, $n$ tane $\TMstroke$ dizisinde
ancak ve ancak $n$ çift ise; $M'$ ise $n$ tane $A$ dizisinde ancak ve ancak
$n$ çift ise durur. Yalnızca $\TMstroke$ simgesini $A$, $q_0$ durumunu $s$
ve $q_1$ durumunu $h$ olarak yeniden adlandırdık. Bu örnek genelleştirilebilir:
Biri ötekinden simge ve durumların bu tür bir yeniden adlandırılmasıyla elde
ediliyorsa Turing makinelerini aynı sayabiliriz. Aslında bir Turing
makinesinin simge ve durumlarını doğrudan pozitif tam sayılar olarak
düşünebiliriz: $\sigma_0$ yerine $1$, $\sigma_1$ yerine $2$ vb.;
$\TMendtape$ simgesi $1$, $\TMblank$ simgesi $2$ vb. olsun. Böylece
\emph{Çift Sayı} makinesi \olref{fig:standard-even} içinde gösterilen
makineye dönüşür.
\begin{figure}
\[\begin{tikzpicture}[->,>=stealth',shorten >=1pt,auto,node distance=2.8cm,
  semithick]
\tikzstyle{every state}=[fill=none,draw=black,text=black]

\node[initial,state]         (A)              {$1$};
\node[state]         (B) [right of=A] {$2$};

\path (A) edge [bend left] node {\TMtrans{3}{3}{\TMright}} (B)
(B) edge [loop above] node {\TMtrans{2}{2}{\TMright}} (B)
edge [bend left] node {\TMtrans{3}{3}{\TMright}} (A);
\end{tikzpicture}
\]
\caption{Standart bir \emph{Çift Sayı} makinesi}
\ollabel{fig:standard-even}
\end{figure}
Durumları ve simgeleri pozitif tam sayılar olan bir Turing makinesine
\emph{standart} makine diyebilir ve bundan sonra yalnızca standart makineleri
ele alabiliriz.\footnote{``Standart makine'' terminolojisi standart değildir.}

Turing makineleri kümesinin !!{enumerable} olduğunu göstermek istiyorduk;
yukarıdaki değerlendirmeler göz önüne alındığında standart Turing makineleri
kümesinin !!{enumerable} olduğunu göstermek yeterlidir. Standart bir
$M=\tuple{Q,\Sigma,q_0,\delta}$ Turing makinesi verildiğini varsayalım. Onu
pozitif tam sayılardan oluşan sonlu bir dizgeyle nasıl betimleyebiliriz? Önce
durumların sayısını, durumların kendilerini, simgelerin sayısını, simgelerin
kendilerini ve başlangıç durumunu listeleriz. (Unutmayın, $M$ standart bir
makine olduğundan bunların hepsi pozitif tam sayıdır.) Peki $\delta$ ne olacak?
Olası argümanlar kümesi, yani $\tuple{q,\sigma}$ çiftleri, $Q$ ile $\Sigma$
sonlu olduğu için sonludur. Dolayısıyla $\delta$'daki bilgi, burada
$\delta(q,\sigma)=\tuple{q',\sigma',D}$ ve $d$ yön~$D$'yi kodlayan bir sayı
(örneğin $\TMleft$ için $1$, $\TMright$ için $2$, $\TMstay$ için $3$) olmak
üzere $\tuple{q,\sigma,q',\sigma',d}$ biçimindeki bütün $5$-lilerin sonlu
listesidir.

Böylece her standart Turing makinesi pozitif tam sayılardan oluşan sonlu bir
listeyle, yani $s_M\in(\PosInt)^*$ dizisiyle betimlenebilir. Örneğin standart
\emph{Çift Sayı} makinesi şu diziyle kodlanır:
\[
2, \underbrace{1, 2}_Q, 3, \overbrace{1, 2, 3}^\Sigma, 1, \underbrace{1, 3, 2, 3, 2}_{\delta(1,3) = \tuple{2,3,R}}, 
\overbrace{2, 2, 2, 2, 2}^{\delta(2,2) = \tuple{2,2,R}},
\underbrace{2, 3, 1, 3, 2}_{\delta(2,3) = \tuple{1,3,R}}.
\]
\end{explain}

\begin{thm}
$\Nat$'ten~$\Nat$'e Turing hesaplanabilir olmayan fonksiyonlar vardır.
\end{thm}

\begin{proof}
Pozitif tam sayıların sonlu dizileri kümesi~$(\PosInt)^*$'ın !!{enumerable}
olduğunu biliyoruz (\cref{sfr:siz:zigzag:prob:posint-star}). Dolayısıyla
standart Turing makinelerinin betimleri kümesi, $(\PosInt)^*$'ın bir alt
kümesi olarak, kendisi de sayılabilirdir. $\Nat$'ten~$\Nat$'e her Turing
hesaplanabilir fonksiyon bir (aslında birçok) Turing makinesince hesaplanır.
Durum ve simgelerini pozitif tam sayılarla yeniden adlandırarak (özellikle
$\TMendtape$ simgesini $1$, $\TMblank$ simgesini $2$ ve $\TMstroke$ simgesini
$3$ yaparak) her Turing hesaplanabilir fonksiyonun standart bir Turing
makinesince hesaplandığını görebiliriz. Bu, $\Nat$'ten~$\Nat$'e bütün Turing
hesaplanabilir fonksiyonlar kümesinin de sayılabilir olduğu anlamına gelir.

Öte yandan $\Nat$'ten~$\Nat$'e bütün fonksiyonlar kümesi !!{enumerable}
değildir (\cref{sfr:siz:red:prob:nat-nat}). Bütün fonksiyonlar bir Turing
makinesince hesaplanabilir olsaydı onları hesaplayan bütün Turing makinesi
betimlerini listeleyerek bütün fonksiyonlar kümesini numaralandırabilirdik.
Dolayısıyla Turing hesaplanabilir olmayan bazı fonksiyonlar vardır.
\end{proof}

\begin{prob}
  Durumların ve alfabe simgelerinin açıkça listelenmesini gerektirmeyen bir
  Turing makinesi betimleme yolu düşünebilir misiniz? Kendi ``standart''
  makine kavramınızı tanımlayabilirsiniz; ama her Turing makinesinin yeni
  anlamınızdaki bir ``standart'' makine tarafından nasıl hesaplanabildiği
  tarafından nasıl benzetilebildiği hakkında bir şey söyleyin.
\end{prob}

\end{document}
